\section{Materiales y métodos}

\subsection{Base de datos}
Se usó la HMC Sleep Staging Database v1.1~\cite{alvarez2022}, con 151 PSG de adultos registradas a \SI{256}{\hertz}, con cuatro canales de EEG (F4, C4, O2 y C3, referidos a la mastoide contralateral), dos de EOG, EMG de mentón y ECG, anotados por técnicos según la AASM. Se eligieron al azar 70 registros con una semilla fija y se separaron por paciente en 56 de entrenamiento y 14 de prueba. Todos los parámetros se eligieron con los pacientes de entrenamiento. La clasificación y la evaluación se restringieron a las épocas entre las marcas de luces apagadas y encendidas, porque la vigilia fuera de ese período es fácil de clasificar y corre hacia W la mediana de la noche con que se normalizan algunos descriptores (Sección~\ref{sec:clasificacion}). El código está en Python y los archivos del repositorio se describen en el Anexo~\ref{sec:archivos}.

\subsection{Preprocesamiento}
\textbf{Exploratorio y filtrado.} En el espectro crudo de cada canal se midió cuánto sobresalen 50 y \SI{60}{\hertz} respecto de sus vecinas, y del encabezado EDF se leyó el rango físico (\SI{\pm800}{\micro\volt} en el EEG) y los filtros del equipo, un pasa-altos de \SI{0,2}{\hertz} y un pasa-bajos de 35 o \SI{128}{\hertz} según el registro. El EEG y el EOG se filtraron con un pasa-altos de \SI{0,2}{\hertz}, un pasa-bajos de \SI{45}{\hertz} y un notch de \SI{50}{\hertz}, y el EMG con un pasa-altos de \SI{10}{\hertz} y el notch. El pasa-altos bajó respecto de los \SI{0,5}{\hertz} del anteproyecto para no recortar las ondas lentas. Se comparó un IIR Butterworth de orden 4 aplicado hacia adelante y hacia atrás con un FIR por ventana de Hamming, y se adoptó el IIR porque da la misma forma de onda con orden 4 en lugar de 4225 coeficientes.

\textbf{Artefacto ocular.} Con E1$-$E2 como referencia se compararon NLMS y el filtro de Wiener, ambos filtros óptimos. Como la referencia también contiene ondas lentas cerebrales, se agregó una compuerta para que el filtro se ajuste y se aplique solo donde el valor eficaz del EOG en \SI{2}{\second} supera el doble de su mediana de la noche. Las variantes se compararon en F4 con la potencia delta y las ondas lentas detectadas por estadio, que deberían bajar en W y REM y no cambiar en N3.

\textbf{Latidos y artefacto cardíaco.} Los latidos se detectaron con Pan-Tompkins~\cite{pan1985} (pasa-banda de 5 a \SI{15}{\hertz}, derivada, cuadrado, integración en \SI{150}{\milli\second} y umbrales adaptativos), con la polaridad decidida una vez por noche y los umbrales reiniciados tras \SI{5}{\second} sin latidos. Para ver si el artefacto cardíaco era relevante, como pidió la cátedra, se promedió el EEG entre $-0{,}2$ y \SI{0,6}{\second} alrededor de cada onda R, donde el EEG se cancela y el artefacto se conserva, con un control en instantes al azar. Para removerlo, en cada latido se restó el promedio de los 30 latidos anteriores, sin las ventanas de más de \SI{500}{\micro\volt} pico a pico.

\textbf{Segmentación.} Siguiendo la devolución, se analizaron épocas fijas de \SI{30}{\second}. Para marcar las que mezclan dos estados se ajustó un modelo AR de orden 8 (Yule-Walker) en los primeros \SI{5}{\second} de cada época de C4 y se aplicó a los cinco tramos siguientes. La época se marcó como no estacionaria si el error de predicción de algún tramo superaba 4 veces al del primero. La marca no interviene en la clasificación.

\subsection{Descriptores}
\textbf{Espectro multitaper.} El periodograma de una época tiene un error del orden de su propio valor. Welch~\cite{welch1967} lo reduce promediando tramos más cortos, mientras que el multitaper~\cite{thomson1982,prerau2017} usa la época completa multiplicada por $K$ ventanas de Slepian ortogonales y promedia los $K$ espectros, casi independientes, con lo que la varianza baja unas $K$ veces. Con $NW = 15$ la resolución es de \SI{1}{\hertz} y se usaron $K = 2NW - 1 = 29$ ventanas. Para H1, Welch se calculó con ventanas de Hann de \SI{4}{\second}, también de \SI{1}{\hertz} de resolución. Por canal se calcularon la potencia de $\delta$, $\theta$, $\alpha$, $\sigma$ (\SIrange{11}{16}{\hertz}) y $\beta$ (\SIrange{13}{30}{\hertz}) en logaritmo y como fracción de la potencia entre 0,5 y \SI{30}{\hertz}, los cocientes $\delta/\theta$, $\delta/\beta$ y $\theta/\alpha$ y la pendiente del espectro en escala log-log entre 1 y \SI{40}{\hertz}~\cite{donoghue2020}, 56 descriptores en total.

\textbf{Husos, ondas lentas y Woody.} Los husos se detectaron en C4~\cite{molle2002} filtrando entre 11 y \SI{16}{\hertz}, con la envolvente de Hilbert suavizada en \SI{0,2}{\second} y un umbral de la media más 1,5 desvíos de la envolvente de la noche, durante 0,5 a \SI{2}{\second}. Las ondas lentas se detectaron en F4 con el criterio visto en clase~\cite{piantoni2013}, filtrando entre 0,2 y \SI{4}{\hertz} y aceptando las ondas con pico negativo menor a \SI{-60}{\micro\volt} y al menos \SI{75}{\micro\volt} pico a pico. Para estimar la forma típica de las ondas de cada paciente se recortó cada una \SI{\pm1}{\second} alrededor de su pico. Como el ruido corre el pico detectado, el promedio ordinario mezcla ondas desalineadas. El método de Woody~\cite{woody1967} parte de ese promedio, busca para cada onda el corrimiento de hasta \SI{0,25}{\second} que maximiza su correlación cruzada con la plantilla, recalcula la plantilla y repite, con un máximo de 20 iteraciones. Por época se guardaron la cantidad y amplitud de husos y de ondas, y el parecido medio de las ondas con la plantilla.

\textbf{CSP.} Los patrones espaciales comunes~\cite{ramoser2000} buscan una combinación de los cuatro canales, $s = \mathbf{w}^\top\mathbf{x}$, con mucha potencia en un estadio y poca en otro. Se aplicaron al par N1 y REM, el que más se confunde. Con las covarianzas medias de entrenamiento de cada estadio, normalizadas por su traza, se resolvió $\mathbf{C}_{N1}\mathbf{w} = \lambda(\mathbf{C}_{N1}+\mathbf{C}_{R})\mathbf{w}$, se conservaron los dos autovectores de cada extremo y el descriptor fue $\log(\mathbf{w}^\top\mathbf{C}\,\mathbf{w})$ de cada época.

\textbf{EOG, EMG y HRV.} Del EOG se usó la energía de E1$-$E2 por época, donde los movimientos oculares se refuerzan y la actividad cerebral se cancela, y del EMG su valor eficaz, ambos en logaritmo. Para la HRV se descartaron los latidos con correlación menor a 0,8 con el latido típico de la noche (la mediana de todos), y después los intervalos RR fuera de 0,3 a \SI{2}{\second} o a más de un 20\,\% de la mediana de sus vecinos. Como la serie queda con muestras irregulares, el espectro se estimó con el periodograma de Lomb en una ventana de \SI{5}{\minute} centrada en cada época, y se guardaron el logaritmo de HF y de LF/HF~\cite{taskforce1996}.

\subsection{Clasificación}
\label{sec:clasificacion}
El anteproyecto proponía un LDA y un HMM, pero la devolución aclaró que la clasificación no puede usar inteligencia artificial, y ambos se entrenan con épocas etiquetadas. Se los reemplazó por reglas construidas a partir del manual AASM~\cite{berry2017} y un suavizado con reglas fijas, lo que además deja cada decisión explicada por una regla concreta.

Las potencias relativas se usaron como fracción y los husos y ondas como conteo por época, magnitudes que no dependen de la amplitud del registro. El EMG, el EOG y LF/HF sí dependen del paciente y se expresaron como múltiplo de la mediana de su noche, $10^{\,x-\operatorname{mediana}(x)}$. Una primera versión con todos los descriptores en percentiles de la noche dio $\kappa = 0{,}245$, porque un umbral en percentiles fija qué fracción de cada noche cumple la regla sin importar cuánto duerme el paciente en cada estadio. Las reglas se evalúan en el orden de la Tabla~\ref{tab:reglas} y gana la primera que se cumple. Una época que cae en N1 por descarte después de una N2 se asigna a N2, según el criterio de continuidad de la AASM, y eso se propaga hasta que otra regla corta la secuencia.

\begin{table}[htbp]
  \centering\small
  \caption{Reglas en orden de evaluación. EMG, EOG y LF/HF en veces la mediana de la noche.}
  \label{tab:reglas}
  \begin{tabularx}{\textwidth}{clXX}
    \toprule
    & Estadio & Condición & Fundamento \\
    \midrule
    1 & W   & EMG $\geq 1{,}3$ y ($\alpha_\text{rel}$(O2) $\geq 0{,}15$ o EOG $\geq 1{,}5$) & Tono muscular, y ritmo $\alpha$ o parpadeos \\
    2 & N3  & $\delta_\text{rel}$(F4) $\geq 0{,}82$ y ondas lentas $\geq 2$ & Ondas lentas de gran amplitud \\
    3 & REM & EMG $\leq 1{,}0$, (EOG $\geq 1{,}6$ o LF/HF $\geq 1{,}5$) y ningún huso & Atonía, movimientos oculares y predominio simpático \\
    4 & N2  & al menos un huso & Husos de sueño \\
    5 & N1  & ninguna de las anteriores & Estadio de transición \\
    \bottomrule
  \end{tabularx}
\end{table}

Cada umbral se fijó mirando, con los pacientes de entrenamiento, el histograma del descriptor en el estadio que la regla detecta y en los que debe separar, cerca del punto donde se cruzan (Fig.~\ref{fig:umbrales} del anexo). Donde el descriptor separa poco y es una alternativa dentro de su regla ($\alpha$ en W y LF/HF en REM) se dejó un umbral más exigente que el cruce. No se buscaron combinaciones de umbrales que maximicen el acuerdo. Como desempate entre N1 y REM se evaluó el puntaje CSP, que cambia una época de uno a otro estadio solo si queda a más de un rango intercuartil típico del punto medio entre las medianas de ambos. Por último, una época distinta de sus dos vecinas iguales toma el estadio de ellas y un tramo de REM de una sola época toma el estadio anterior. Como alternativa se evaluó un filtro de moda de 5 épocas.

\textbf{Interfaz gráfica.} Es una aplicación de escritorio con Tkinter (\texttt{codigo/app/interfaz.py}) que llama a las mismas funciones del pipeline. Se elige el registro y la época, se configura el suavizado y el desempate CSP y se procesa el registro en segundo plano, con un registro de avance por etapa. Arriba, el hipnograma de los técnicos permite ir a cualquier época con un click. Las pestañas muestran los ocho canales crudos, cada etapa del preprocesamiento sobre el canal elegido, los latidos con el descarte por forma, los intervalos RR y el periodograma de Lomb, y la clasificación con ambos hipnogramas, la matriz de confusión, el $\kappa$ del registro y la regla y los descriptores de cada época (Fig.~\ref{fig:interfaz} del anexo).

\subsection{Evaluación}
Sobre los pacientes de prueba se calculó el $\kappa$ de Cohen, $\kappa = (p_o - p_e)/(1 - p_e)$, donde $p_o$ es la fracción de épocas en que el algoritmo coincide con los técnicos y $p_e$ la esperada por azar, y el F1 de cada estadio, que no se infla con los estadios frecuentes. Cada hipótesis se verificó comparando dos variantes sobre los mismos pacientes de prueba con la prueba de Wilcoxon de rangos con signo. H1 midió la diferencia absoluta media del logaritmo de la potencia entre épocas consecutivas de un mismo estadio, por banda y frecuencia por frecuencia entre 0,5 y \SI{30}{\hertz}. H2 insertó 40 ondas lentas conocidas (coseno de \SI{0,8}{\hertz} con envolvente gaussiana y \SI{80}{\micro\volt} de pico) con corrimientos al azar de hasta \SI{0,25}{\second} en EEG real de vigilia y midió el error cuadrático medio de la forma recuperada, relativo al valor eficaz de la onda. H3 comparó el F1 de N1 y de REM con y sin el desempate CSP. H4 midió el residuo cardíaco, como el máximo absoluto del promedio del EEG alrededor de los latidos promediado en los cuatro canales, y el $\kappa$ con y sin remoción de artefactos y HRV.
